% Báo cáo SCCG — dựa theo cấu trúc VDT2026_MiniProject_report
\documentclass[11pt,a4paper,oneside]{report}

\usepackage{fontspec}
\usepackage[vietnamese,english]{babel}
\usepackage{geometry}
\geometry{margin=2.5cm}
\usepackage{setspace}
\onehalfspacing
\usepackage{graphicx}
\usepackage{booktabs}
\usepackage{longtable}
\usepackage{array}
\usepackage{float}
\usepackage{hyperref}
\usepackage{xurl}
\usepackage{enumitem}
\usepackage{tikz}
\usetikzlibrary{arrows.meta,positioning,shapes.geometric}
\usepackage{csquotes}
\usepackage[backend=biber,style=numeric,sorting=none,maxbibnames=99]{biblatex}
\addbibresource{references.bib}
\hypersetup{colorlinks=true,linkcolor=blue,urlcolor=blue,citecolor=blue}
\setlength{\parskip}{0.35em}
\setlength{\parindent}{1.2em}
\renewcommand{\chaptername}{Chương}

\newcommand{\reportauthor}{Bùi Đình Bảo}
\newcommand{\reportprogram}{Chương trình Viettel Digital Talent 2026}
\newcommand{\reporttitle}{SCCG: Đề xuất kiểm tra tuân thủ OSS trong quy trình phát triển phần mềm}
\newcommand{\reportdate}{Tháng 10 năm 2026}
\newcommand{\reportstatus}{Báo cáo đề xuất và thiết kế pilot}

\begin{document}
\begin{titlepage}
\centering
\vspace*{1.2cm}
{\Large\bfseries TẬP ĐOÀN CÔNG NGHIỆP -- VIỄN THÔNG QUÂN ĐỘI\par}
\vspace{2.8cm}
{\Large\bfseries BÁO CÁO SÁNG KIẾN VDT\par}
\vspace{1cm}
{\huge\bfseries \reporttitle\par}
\vspace{1.5cm}
{\Large \reportauthor\par}
\vspace{0.5cm}
{\large \reportprogram\par}
\vfill
{\large Trạng thái: \reportstatus\par}
\vspace{0.4cm}
{\large \reportdate\par}
\vspace{1cm}
\end{titlepage}

\pagenumbering{roman}
\chapter*{Lời cảm ơn}
Báo cáo được hình thành từ trải nghiệm cá nhân trong quá trình thử nghiệm các thư viện, mô hình và công cụ mã nguồn mở phục vụ pipeline AI OCR. Tôi xin cảm ơn chương trình Viettel Digital Talent đã tạo môi trường để học tập, thực hành kỹ thuật và đề xuất những cải tiến có thể kiểm chứng. Các nội dung về quy trình, chính sách và khả năng áp dụng trong tổ chức cần tiếp tục được trao đổi với các nhóm kỹ thuật, DevSecOps và bộ phận phụ trách pháp chế/tuân thủ trước khi triển khai thử nghiệm.

\chapter*{Lời mở đầu}
Phần mềm hiện đại thường được xây dựng từ nhiều thành phần: thư viện, package, công cụ build, container, mô hình AI, trọng số và dataset. Việc lựa chọn một thành phần không chỉ phụ thuộc vào chức năng hay hiệu năng; nhóm phát triển còn cần hiểu điều kiện sử dụng, nguồn gốc và bằng chứng liên quan đến phiên bản được tích hợp. Những thông tin này có thể nằm ở nhiều nguồn khác nhau và mức độ đầy đủ không đồng nhất.

Từ trải nghiệm thử nghiệm OSS trong pipeline AI OCR, tôi đặt câu hỏi liệu việc kiểm tra giấy phép, dependency và nguồn bằng chứng có thể được thực hiện sớm, nhất quán và có thể truy vết hơn hay không. Câu hỏi này là giả thuyết cá nhân, không phải kết luận rằng tổ chức hiện chưa có quy trình tương đương hoặc đã xảy ra vi phạm.

Báo cáo đề xuất SCCG (Source Code Compliance Gate) như một bước hỗ trợ sàng lọc. Trọng tâm của giai đoạn đầu không phải xây dựng ngay một nền tảng CI/CD toàn diện, mà là kiểm chứng một luồng hẹp: thu thập thông tin của thành phần OSS, đánh giá chất lượng bằng chứng, đối chiếu với chính sách do đơn vị có thẩm quyền cung cấp và xuất báo cáo giải thích được. Các quyết định mơ hồ cần được chuyển cho con người xem xét; công cụ không thay thế tư vấn pháp lý.

\tableofcontents
\listoftables
\clearpage
\pagenumbering{arabic}

\chapter{Giới thiệu}

\section{Đặt vấn đề}
Trong quá trình phát triển phần mềm, nhóm kỹ thuật thường sử dụng thư viện, package, công cụ build và tài nguyên từ bên thứ ba để rút ngắn thời gian phát triển. Với các hệ thống AI, phạm vi tài nguyên có thể rộng hơn, gồm mã nguồn mô hình, trọng số, dataset và công cụ xử lý dữ liệu. Mỗi loại tài nguyên có thể đi kèm điều kiện sử dụng khác nhau; không nên suy luận rằng giấy phép được ghi ở một nơi tự động áp dụng cho toàn bộ tài nguyên liên quan.

Thông tin về giấy phép, phiên bản và nguồn gốc có thể phân tán giữa repository, package registry, model card, tệp manifest, lockfile và tài liệu phát hành. Nếu việc kiểm tra chưa được tích hợp vào thời điểm phù hợp, người phát triển có thể phải tra cứu lặp lại hoặc xử lý lại khi thành phần đã được tích hợp sâu. Tuy nhiên, mức độ của vấn đề trong quy trình hiện tại cần được khảo sát; báo cáo không giả định rằng tổ chức chưa có công cụ hay quy trình tương đương.

\section{Nguồn gốc ý tưởng}
Ý tưởng bắt đầu từ trải nghiệm cá nhân khi thử nghiệm nhiều mô hình, repository và công cụ OSS cho pipeline AI OCR. Trong quá trình so sánh chất lượng, tốc độ và khả năng tích hợp, trọng tâm dễ nghiêng về khía cạnh kỹ thuật trong khi điều kiện sử dụng và nguồn bằng chứng cần được kiểm tra riêng.

Từ đó, câu hỏi được mở rộng thành: làm thế nào để hỗ trợ nhóm phát triển nhận biết sớm những thành phần cần rà soát, phân biệt trường hợp thiếu thông tin với trường hợp không phù hợp chính sách, và lưu lại bằng chứng phục vụ kiểm tra sau này?

\section{Mục tiêu báo cáo}
Báo cáo có bốn mục tiêu:
\begin{enumerate}[label=\arabic*.]
    \item Mô tả vấn đề và các giả thuyết cần xác minh với nhóm kỹ thuật, DevSecOps và bộ phận phụ trách pháp chế/tuân thủ.
    \item Đề xuất phạm vi MVP của SCCG, ưu tiên thu thập bằng chứng, phân tích dependency trong phạm vi hỗ trợ và giải thích kết quả.
    \item Xác định kiến trúc, yêu cầu chức năng, rủi ro và quy trình pilot có thể kiểm chứng.
    \item Đề xuất phương pháp đánh giá dựa trên baseline, kết quả rà soát thủ công và số liệu thực nghiệm thay vì khẳng định trước hiệu quả.
\end{enumerate}

\section{Câu hỏi cần trả lời}
\begin{enumerate}[label=\textbf{RQ\arabic*:}]
    \item Có khoảng trống nào trong quy trình hiện tại liên quan đến kiểm tra giấy phép, nguồn bằng chứng và dependency OSS?
    \item Có thể thu thập được thông tin đủ tin cậy để gắn với thành phần và phiên bản cụ thể đến mức nào?
    \item SCCG có thể phân biệt kết quả chính sách với chất lượng bằng chứng và đưa ra lý do dễ kiểm tra không?
    \item So với rà soát thủ công và công cụ hiện có, SCCG có cải thiện độ bao phủ hoặc thời gian xử lý mà không tạo quá nhiều cảnh báo sai không?
\end{enumerate}

\section{Phạm vi và giới hạn}
MVP tập trung vào một hệ sinh thái package/repository được chọn sau khảo sát, đầu vào là URL hoặc danh sách thành phần và đầu ra là báo cáo có nguồn bằng chứng. Dependency trực tiếp và gián tiếp chỉ được phân tích trong phạm vi mà công cụ có thể tái dựng đáng tin cậy.

MVP chưa nhằm tự kết luận vi phạm pháp luật, chứng minh nguồn gốc mọi dòng mã, phát hiện chắc chắn mã do AI sinh ra, hay bao phủ toàn bộ mô hình, dataset, binary và container. Không được gửi mã nguồn hoặc dữ liệu nội bộ tới dịch vụ bên ngoài khi chưa được phê duyệt. Việc tích hợp CI là bước sau pilot và ban đầu chỉ chạy ở chế độ cảnh báo.

\chapter{Nội dung và thiết kế giải pháp}

\section{Khảo sát vấn đề và giả thuyết}
SCCG xuất phát từ giả thuyết rằng một luồng kiểm tra có cấu trúc có thể giảm thao tác tra cứu lặp lại và giúp lưu bằng chứng xử lý. Trước khi phát triển, cần khảo sát quy trình hiện tại, công cụ đã được sử dụng, loại dự án cần ưu tiên, chủ sở hữu chính sách và các yêu cầu bảo mật.

Các câu hỏi khảo sát gồm: (i) kiểm tra license/OSS hiện diễn ra ở giai đoạn nào; (ii) có công cụ SBOM hoặc license scanning nào đang vận hành; (iii) cách xử lý trường hợp metadata thiếu hoặc mâu thuẫn; (iv) ai phê duyệt ngoại lệ; (v) dữ liệu nào được phép đưa vào công cụ; và (vi) chỉ số nào có thể dùng làm baseline.

\section{Các khái niệm cần phân biệt}
\subsection{Thành phần, phiên bản và nguồn bằng chứng}
Một kết quả cần được gắn với định danh thành phần và phiên bản/commit cụ thể. Nguồn bằng chứng có thể gồm metadata từ registry, tệp LICENSE, thông tin phát hành, manifest, lockfile hoặc model card. Sự hiện diện của trường license không tự chứng minh rằng thông tin đó đầy đủ, chính xác hay áp dụng cho mọi thành phần đi kèm.

\subsection{Kết quả chính sách và chất lượng bằng chứng}
Hai khái niệm này phải được biểu diễn riêng. Kết quả chính sách phản ánh việc bằng chứng đã thu thập được có phù hợp với quy tắc đã được phê duyệt hay không; chất lượng bằng chứng phản ánh mức độ đầy đủ, truy xuất được và nhất quán của thông tin. Thiếu bằng chứng không đồng nghĩa với vi phạm.

\begin{table}[H]
\centering
\caption{Phân biệt kết quả chính sách và trạng thái bằng chứng}
\label{tab:evidence-policy}
\begin{tabular}{p{0.24\textwidth}p{0.31\textwidth}p{0.34\textwidth}}
\toprule
\textbf{Chiều đánh giá} & \textbf{Giá trị đề xuất} & \textbf{Ý nghĩa} \\
\midrule
Kết quả chính sách & Đạt / Cần xem xét / Không đạt & So sánh bằng chứng với quy tắc được phê duyệt \\
Chất lượng bằng chứng & Đầy đủ / Một phần / Mâu thuẫn hoặc chưa xác minh & Cho biết mức độ tin cậy và khả năng truy vết \\
Hành động tiếp theo & Chấp nhận / Rà soát / Thay thế / Xin ngoại lệ & Hướng xử lý, không phải kết luận pháp lý tự động \\
\bottomrule
\end{tabular}
\end{table}

\section{Yêu cầu chức năng}
Các yêu cầu dưới đây là mục tiêu thiết kế của MVP, chưa phải tuyên bố về tính năng đã được triển khai.

\begin{longtable}{p{0.13\textwidth}p{0.40\textwidth}p{0.37\textwidth}}
\caption{Yêu cầu chức năng mức MVP}\label{tab:fr}\\
\toprule
\textbf{Mã} & \textbf{Yêu cầu} & \textbf{Tiêu chí kiểm chứng} \\
\midrule
\endfirsthead
\toprule
\textbf{Mã} & \textbf{Yêu cầu} & \textbf{Tiêu chí kiểm chứng} \\
\midrule
\endhead
FR-01 & Nhận URL repository/package/model hoặc danh sách thành phần & Đầu vào được kiểm tra và định danh thành phần được ghi lại \\
FR-02 & Thu thập phiên bản, metadata license, nguồn tham khảo và thời điểm kiểm tra & Mỗi trường có nguồn hoặc được đánh dấu chưa có bằng chứng \\
FR-03 & Thu thập dependency trực tiếp/gián tiếp trong hệ sinh thái được chọn & Báo cáo có thể truy vết đường dẫn dependency khi dữ liệu hỗ trợ \\
FR-04 & Đối chiếu với chính sách có phiên bản do chủ sở hữu chính sách cung cấp & Kết quả chỉ dựa trên quy tắc đã cấu hình và được ghi phiên bản \\
FR-05 & Phân biệt đạt, cần xem xét và không đạt & Thiếu hoặc mâu thuẫn dữ liệu không bị tự động coi là vi phạm \\
FR-06 & Xuất báo cáo đọc được và JSON có cấu trúc & Có thành phần, phiên bản, kết quả, lý do, nguồn và ngày kiểm tra \\
FR-07 & Ghi nhận ngoại lệ có người phê duyệt, lý do và phạm vi & Ngoại lệ có dấu vết kiểm toán và không bị bỏ qua âm thầm \\
FR-08 & So sánh kết quả tự động với tập mẫu đã rà soát thủ công & Có bảng lỗi, trường hợp chưa kết luận và số liệu đánh giá \\
\bottomrule
\end{longtable}

\section{Yêu cầu phi chức năng}
\begin{itemize}
    \item \textbf{Khả năng truy vết:} kết quả gắn với thành phần/phiên bản, nguồn bằng chứng, thời điểm kiểm tra và phiên bản chính sách.
    \item \textbf{Tính tái lập:} cùng đầu vào, bằng chứng và phiên bản chính sách phải tạo ra kết quả nhất quán, ngoại trừ dữ liệu bên ngoài đã thay đổi.
    \item \textbf{Bảo mật:} không chuyển source code hoặc metadata nội bộ ra ngoài môi trường được phê duyệt.
    \item \textbf{Khả năng giải thích:} mỗi cảnh báo chỉ ra quy tắc, bằng chứng và hành động tiếp theo.
    \item \textbf{Khả năng tích hợp:} ưu tiên định dạng đầu ra có thể lưu cùng PR và đọc được bởi CI ở giai đoạn sau.
    \item \textbf{Tính an toàn vận hành:} pilot mặc định cảnh báo, không chặn merge.
\end{itemize}

\section{Kiến trúc mức cao}
Kiến trúc đề xuất tách quá trình thu thập dữ liệu khỏi việc áp dụng chính sách. Điều này cho phép kiểm tra nguồn bằng chứng độc lập, thay đổi policy version có kiểm soát và giải thích quyết định sau này.

\begin{figure}[H]
\centering
\begin{tikzpicture}[
  node distance=0.65cm,
  box/.style={draw,rounded corners,align=center,text width=0.78\textwidth,minimum height=0.8cm},
  arrow/.style={-{Stealth[length=2mm]},thick}
]
\node[box] (input) {Đầu vào: repository / package / model URL / manifest};
\node[box,below=of input] (discover) {Khám phá thành phần và dependency trong phạm vi hỗ trợ};
\node[box,below=of discover] (evidence) {Thu thập và chuẩn hóa bằng chứng: phiên bản, license, URL, thời điểm};
\node[box,below=of evidence] (quality) {Đánh giá chất lượng bằng chứng: đầy đủ / một phần / mâu thuẫn};
\node[box,below=of quality] (policy) {Policy engine có phiên bản: đối chiếu quy tắc đã được phê duyệt};
\node[box,below=of policy] (report) {Báo cáo: kết quả, lý do, đường dẫn dependency, nguồn và hướng xử lý};
\draw[arrow] (input)--(discover);
\draw[arrow] (discover)--(evidence);
\draw[arrow] (evidence)--(quality);
\draw[arrow] (quality)--(policy);
\draw[arrow] (policy)--(report);
\end{tikzpicture}
\caption{Kiến trúc logic đề xuất cho SCCG MVP}
\label{fig:sccg-architecture}
\end{figure}

\section{Mô hình quyết định và luồng ngoại lệ}
Kết quả nên dùng ba trạng thái chính: \textit{Đạt}, \textit{Cần xem xét} và \textit{Không đạt theo quy tắc đã phê duyệt}. Trạng thái bằng chứng được lưu riêng để tránh biến sự thiếu thông tin thành một kết luận pháp lý. Chỉ người có thẩm quyền mới được xác nhận ngoại lệ; ngoại lệ cần lưu người duyệt, lý do, phạm vi, ngày phê duyệt và thời hạn nếu có.

Trong pilot, tất cả kết quả chỉ tạo báo cáo/cảnh báo. Chỉ cân nhắc chặn merge sau khi có baseline, độ chính xác phù hợp, chủ sở hữu chính sách, quy trình khiếu nại/ngoại lệ và sự đồng ý của các bên liên quan.

\section{Khảo sát công cụ liên quan}
Các công cụ hiện có cần được đánh giá trước khi tự xây dựng toàn bộ pipeline. SPDX cung cấp chuẩn định danh và biểu diễn thông tin license/SBOM; ScanCode Toolkit có thể hỗ trợ quét thông tin license/copyright; Syft tạo SBOM cho nhiều loại artifact; các công cụ như ORT hỗ trợ các bước phân tích tuân thủ license trong quy trình phần mềm \cite{spdx,scancode,syft,ort}. Đây là các ứng viên khảo sát, không phải khẳng định rằng SCCG cần thay thế chúng hoặc đã có tính năng vượt trội.

Việc lựa chọn baseline cần căn cứ vào hệ sinh thái mục tiêu, khả năng chạy nội bộ, độ bao phủ, khả năng xuất bằng chứng, chi phí tích hợp và giấy phép của chính công cụ. Kết quả khảo sát phải được ghi vào pilot trước khi chốt phạm vi phát triển.

\chapter{Kế hoạch pilot và đánh giá}

\section{Trạng thái hiện tại}
Tại thời điểm lập báo cáo, SCCG là đề xuất cần khảo sát và thử nghiệm; tài liệu này không khẳng định MVP đã được triển khai hoặc đã đạt các chỉ số hiệu quả. Các kết quả sẽ chỉ được đưa vào sau khi có tập mẫu, baseline, người rà soát phù hợp và dữ liệu đo có thể kiểm tra lại.

\section{Phương pháp thử nghiệm}
Đề xuất triển khai pilot theo năm giai đoạn:
\begin{enumerate}[label=\arabic*.]
    \item \textbf{Khảo sát quy trình:} trao đổi với nhóm kỹ thuật, DevSecOps và pháp chế/tuân thủ; xác định công cụ hiện có, chính sách được áp dụng, hệ sinh thái mục tiêu và giới hạn dữ liệu.
    \item \textbf{Xây dựng tập mẫu:} chọn khoảng 20--30 thành phần nếu được phép và đủ đa dạng; mỗi mẫu được rà soát thủ công, ghi nguồn bằng chứng và lý do phân loại. Quy mô này chỉ phục vụ thăm dò, không đại diện cho mọi hệ sinh thái.
    \item \textbf{Chạy baseline:} đánh giá một hoặc nhiều công cụ hiện có và quy trình thủ công trên cùng tập mẫu.
    \item \textbf{Xây dựng MVP hẹp:} triển khai thu thập metadata, phân loại bằng chứng, policy engine có phiên bản và xuất báo cáo/JSON.
    \item \textbf{Đối chiếu và quyết định:} so sánh kết quả, phân tích lỗi và quyết định có nên tích hợp một repository vào CI ở chế độ cảnh báo.
\end{enumerate}

\section{Chỉ số đánh giá}
\begin{longtable}{p{0.24\textwidth}p{0.36\textwidth}p{0.30\textwidth}}
\caption{Chỉ số và cách đo trong pilot}\label{tab:metrics}\\
\toprule
\textbf{Chỉ số} & \textbf{Cách đo} & \textbf{Lưu ý} \\
\midrule
\endfirsthead
\toprule
\textbf{Chỉ số} & \textbf{Cách đo} & \textbf{Lưu ý} \\
\midrule
\endhead
Độ bao phủ bằng chứng & Tỷ lệ thành phần có thông tin license và nguồn truy xuất được & Báo cáo riêng theo hệ sinh thái và loại tài nguyên \\
Độ đúng phân loại & So sánh kết quả với nhãn do người có chuyên môn xác minh & Trường hợp chưa kết luận phải ghi riêng \\
Cảnh báo sai & Số cảnh báo không được xác nhận khi rà soát thủ công & Xác định quy tắc và mẫu số trước khi tính tỷ lệ \\
Bỏ sót & Số trường hợp cần cảnh báo theo rà soát thủ công nhưng công cụ không phát hiện & Không gộp với trường hợp thiếu dữ liệu \\
Thời gian xử lý & So sánh thời gian kiểm tra thủ công với quy trình có công cụ & Dùng cùng tập mẫu và ghi rõ thao tác được tính \\
Khả năng truy vết & Tỷ lệ kết quả có phiên bản, nguồn, lý do và policy version & Là điều kiện quan trọng để tái kiểm tra \\
Chi phí vận hành & Thời gian tích hợp, cập nhật rule và bảo trì & Chưa quy đổi thành tiết kiệm tiền nếu thiếu dữ liệu nền \\
Ảnh hưởng CI & Thời gian tăng thêm và số cảnh báo phát sinh trên PR & Chỉ đo sau khi pilot được phê duyệt \\
\bottomrule
\end{longtable}

Với tập mẫu nhỏ, nên báo cáo cả số lượng tuyệt đối lẫn tỷ lệ và không diễn giải kết quả như bằng chứng thống kê cho toàn tổ chức. Khi có đủ mẫu, có thể báo cáo precision, recall và các trường hợp false positive/false negative theo từng nhóm rủi ro.

\section{Tiêu chí nghiệm thu đề xuất}
\begin{enumerate}[label=\arabic*.]
    \item Mỗi kết quả truy ngược được về thành phần/phiên bản, nguồn bằng chứng, thời điểm kiểm tra và phiên bản chính sách.
    \item Thiếu hoặc mâu thuẫn bằng chứng được phân biệt với kết quả không phù hợp chính sách.
    \item Dependency gián tiếp trong phạm vi hỗ trợ có thể truy vết theo đường dẫn.
    \item Có báo cáo so sánh với baseline, bao gồm cảnh báo sai, bỏ sót, trường hợp chưa kết luận và thời gian xử lý.
    \item Có hướng dẫn rà soát và ngoại lệ; chế độ mặc định không chặn merge.
    \item Có đánh giá bảo mật và xác nhận phạm vi dữ liệu trước khi dùng trong repository thật.
\end{enumerate}

\section{Rủi ro và biện pháp giảm thiểu}
\begin{longtable}{p{0.27\textwidth}p{0.58\textwidth}}
\caption{Rủi ro chính của SCCG pilot}\label{tab:risks}\\
\toprule
\textbf{Rủi ro} & \textbf{Biện pháp giảm thiểu} \\
\midrule
\endfirsthead
\toprule
\textbf{Rủi ro} & \textbf{Biện pháp giảm thiểu} \\
\midrule
\endhead
Metadata thiếu hoặc sai & Lưu nguồn, thời điểm, phiên bản; hiển thị mức độ đầy đủ và chuyển trường hợp mơ hồ cho người rà soát \\
Cảnh báo sai làm gián đoạn phát triển & Bắt đầu ở chế độ báo cáo; đo false positive trước khi cân nhắc enforcement \\
Bỏ sót dependency & So sánh với lockfile/SBOM và rà soát thủ công; công bố hệ sinh thái chưa hỗ trợ \\
Lộ source code hoặc metadata nội bộ & Ưu tiên chạy trong hạ tầng được phê duyệt; không gửi dữ liệu ra ngoài khi chưa có cho phép \\
Chính sách không rõ hoặc lỗi thời & Xác định chủ sở hữu, quản lý phiên bản policy và ngày hiệu lực \\
Trùng lặp công cụ hiện có & Thực hiện khảo sát và baseline trước khi xây; ưu tiên tích hợp công cụ phù hợp \\
Hiểu nhầm kết quả thành tư vấn pháp lý & Ghi rõ giới hạn, lý do và quy trình chuyển cấp cho người có thẩm quyền \\
\bottomrule
\end{longtable}

\section{Đối chiếu tiêu chí đánh giá sáng kiến VDT}
\subsection{Hiệu quả kinh tế}
Giá trị tiềm năng có thể đến từ giảm thời gian tra cứu lặp lại, phát hiện vấn đề sớm hơn và giảm công sức sửa đổi khi thành phần đã được tích hợp. Cần đo thời gian nền và chi phí vận hành trước khi tính hiệu quả. Báo cáo không đưa ra tỷ lệ tiết kiệm chưa được đo.

\subsection{Tính mới}
Quét license và tạo SBOM đã có công cụ chuyên dụng. Vì vậy, bản thân chức năng quét không được coi là tính mới. Cần xác định khoảng trống sau khảo sát; khác biệt có thể nằm ở tích hợp quy trình nội bộ, quản lý bằng chứng, giải thích policy và luồng ngoại lệ, nhưng phải chứng minh bằng so sánh.

\subsection{Tính ứng dụng}
Có thể bắt đầu từ một hệ sinh thái và một repository thử nghiệm, nếu phù hợp với quy trình, hạ tầng và chính sách hiện có. Mức độ sẵn sàng chỉ được đánh giá sau khi xác định được người sở hữu chính sách và điều kiện dữ liệu.

\subsection{Tính hoàn thiện}
Tài liệu thiết kế không thay thế sản phẩm và kết quả pilot. Hồ sơ hoàn thiện cần có MVP chạy được, tập mẫu đã xác minh, báo cáo đo, hướng dẫn vận hành, phân tích lỗi và giới hạn sử dụng. Chỉ đánh giá mức hoàn thiện sau khi các tiêu chí này được kiểm chứng.

\chapter{Kết luận và hướng phát triển}

\section{Tổng kết}
SCCG được đề xuất như một bước hỗ trợ kiểm tra việc sử dụng OSS trong quy trình phát triển phần mềm, bắt đầu từ trải nghiệm cá nhân khi thử nghiệm các thành phần OSS cho AI OCR. Vấn đề cần giải quyết chưa được coi là đã xác minh; bước đầu tiên là khảo sát quy trình và công cụ hiện tại, sau đó xác định khoảng trống thực tế.

Thiết kế MVP ưu tiên thu thập bằng chứng gắn với thành phần/phiên bản, phân biệt chất lượng bằng chứng với kết quả chính sách, truy vết dependency trong phạm vi hỗ trợ và tạo báo cáo giải thích được. Policy engine cần dùng quy tắc đã được phê duyệt, có phiên bản và lưu dấu vết xử lý ngoại lệ. Công cụ không tự kết luận vi phạm pháp luật, và thiếu metadata không được xem là bằng chứng tự động của vi phạm.

\section{Đóng góp dự kiến}
Nếu pilot được chấp thuận và triển khai, các đầu ra kỳ vọng gồm:
\begin{itemize}
    \item Bản mô tả quy trình và khoảng trống đã được xác minh với các bên liên quan.
    \item Bộ mẫu nhỏ có kết quả rà soát thủ công và nguồn bằng chứng.
    \item MVP thu thập thông tin, áp dụng quy tắc và xuất báo cáo có thể truy vết.
    \item Báo cáo so sánh với rà soát thủ công và công cụ hiện có.
    \item Khuyến nghị có căn cứ về việc tích hợp CI, mở rộng phạm vi hay dừng thử nghiệm.
\end{itemize}
Đây là đầu ra dự kiến, không phải kết quả đã hoàn thành.

\section{Hạn chế}
Tại thời điểm báo cáo, chưa có số liệu pilot để khẳng định độ chính xác, mức giảm thời gian, hiệu quả kinh tế hoặc tính mới. Chất lượng metadata khác nhau giữa hệ sinh thái; một số tài nguyên AI có thể không cung cấp đầy đủ thông tin về trọng số/dataset. Dependency gián tiếp chỉ có thể phân tích đáng tin cậy khi có manifest, lockfile hoặc dữ liệu tương ứng. Quy tắc chính sách và quyết định ngoại lệ phải do người có thẩm quyền cung cấp.

\section{Hướng phát triển}
Các bước tiếp theo được đề xuất theo thứ tự:
\begin{enumerate}[label=\arabic*.]
    \item Xác minh nhu cầu, công cụ hiện tại, chính sách và điều kiện bảo mật.
    \item Chọn một hệ sinh thái và xác lập baseline.
    \item Xây tập mẫu đã rà soát thủ công và xác định quy tắc đánh giá.
    \item Triển khai MVP hẹp, ưu tiên bằng chứng và giải thích kết quả.
    \item Đánh giá sai số, độ bao phủ và thời gian; chỉ sau đó thử tích hợp CI ở chế độ cảnh báo.
    \item Chỉ cân nhắc chặn merge hoặc mở rộng sang AI provenance khi đã có bằng chứng, quy trình ngoại lệ và phê duyệt phù hợp.
\end{enumerate}

Kết luận cuối cùng về tính khả thi và giá trị của SCCG cần dựa trên kết quả khảo sát và thử nghiệm, không chỉ dựa trên mô tả kiến trúc.


\printbibliography[title={Tài liệu tham khảo}]
\end{document}
